\section{GRPO 的直觉}

\subsection{为什么是 group relative}

GRPO 的关键不在于名字，而在于它利用了语言模型的自然分组结构：
\begin{itemize}
    \item 同一个 prompt 可以采样出多条 response；
    \item 这些 response 天然构成一个 group；
    \item group 内平均 reward 就是一个很自然的 baseline。
\end{itemize}

这就是 “relative” 的来源。它比较的是“同一 prompt 下，谁比同组其他 response 更好”。

\begin{knowledgebox}{为什么 GRPO 适合语言模型}
对同一个 prompt 采样多条 response，本来就是 LM 训练里最自然的采样单位，因此 baseline 不必再从全局状态空间拟合，而可以在 group 内做相对比较。
\end{knowledgebox}

\begin{table}[H]
\centering
\begin{tabular}{@{}p{2.4cm}p{3.8cm}p{3.6cm}p{3.5cm}@{}}
\toprule
\textbf{方法} & \textbf{信号来源} & \textbf{baseline / critic} & \textbf{特点} \\
\midrule
REINFORCE & 直接用 reward & 无显式 critic & 最朴素，方差大 \\
PPO & reward + value estimate & 需要 critic / value model & 更稳定，但系统更重 \\
GRPO & group 内相对 reward & group mean 近似 baseline & 适合多样本 response 的 LM 场景 \\
\bottomrule
\end{tabular}
\caption{三种常见策略梯度风格方法的关系：GRPO 可以看成是 LM 场景下更轻量的 PPO 变体}
\end{table}

\begin{knowledgebox}{为什么 GRPO 不一定要 critic}
critic 的作用是估计 state value，帮助你减少方差。但在同一个 prompt 下采样多个 response 时，group mean 本身就已经提供了一个很自然的、低成本的 baseline 近似，因此可以先不用单独训练 value network。
\end{knowledgebox}



\begin{warningbox}{group mean 也不是万能的}
group mean 是一个很强的经验 baseline，但它并不等于最优 baseline。若 group 太小，估计会噪声大；若 group 太大，采样和计算成本会上升。课程脚本用的是一个教学上很好理解的折中。
\end{warningbox}

\subsection{和 PPO 的关系}

课程里把 GRPO 描述成一种 \textbf{不需要 critic} 的 PPO 风格简化版。你仍然会看到：
\begin{itemize}
    \item policy ratio；
    \item clipping；
    \item 参考策略或旧策略；
    \item KL regularization。
\end{itemize}

但 group 结构让 baseline 近似得更自然，所以实现上可以更轻。


